\documentclass[12pt]{article}

\usepackage{amsmath}
\usepackage{amssymb}
\usepackage{amsthm}
\usepackage{physics}
\usepackage{ragged2e}
\justifying

\setlength{\parindent}{0pt}
\setlength{\parskip}{12pt}
\usepackage[a4paper, margin=2cm]{geometry}

\pagestyle{empty}

\newtheorem*{theorem}{Theorem}
\newtheorem*{lemma}{Lemma}

\begin{document}
\begin{center}
	{\Large \textbf{CSSA HELL}} \\
	Why would you do this
	\rule{\textwidth}{1pt}
\end{center}

We have the following:
\begin{align*}
	a_A & = -(g+kv) \\
	a_B & = -(g+kv).
\end{align*}
Put $a = \dv{v}{t}$ and integrate.
\begin{align*}
	\ln\abs{g+kv_A} & = -kt + C_1  \\
	\ln\abs{g+kv_B} & = -kt + C_2.
\end{align*}
Set $t=0$ for both, with initial unsigned velocities $u_1$, $u_2$, and take upwards as positive to yield:
\begin{align*}
  C_1 & = \ln\abs{g+ku_1} \\
  C_2 & = \ln\abs{g-ku_2}.
\end{align*}
Considering initial signs, we have:
\begin{align*}
	v_A & = \qty(\dfrac{g+ku_1}{k})e^{-kt} - \dfrac{g}{k}   \\
	v_B & = -\qty(\dfrac{ku_2-g}{k})e^{-kt} - \dfrac{g}{k}.
\end{align*}
Now, integrate for $x_A$, $x_B$:
\begin{align*}
	x_A & = \int v_A \dd{t}                                         \\
	    & = C_3 - \qty(\dfrac{g+ku_1}{k^2})e^{-kt} - \dfrac{gt}{k}  \\
	x_B & = \int v_B \dd{t}                                         \\
	    & = C_4 + \qty(\dfrac{ku_2-g}{k^2})e^{-kt} - \dfrac{gt}{k}.
\end{align*}
Taking initial conditions again, we have:
\begin{align*}
	C_3 & = \dfrac{g+ku_1}{k^2}    \\
	C_4 & = H-\dfrac{ku_2-g}{k^2},
\end{align*}
which gives
\begin{align*}
	x_A & = \qty(1-e^{-kt})\qty(\dfrac{g+ku_1}{k^2}) - \dfrac{gt}{k}      \\
	x_B & = H - \qty(1-e^{-kt})\qty(\dfrac{ku_2-g}{k^2}) - \dfrac{gt}{k}.
\end{align*}
\newpage
Suppose $x_A(T)$ is maximised. Then, $x_B(T) - x_A(T) > 0$.
Also,
\[T = \dfrac{1}{k}\ln\qty(\dfrac{g+ku_1}{g}),\]
so
\[  H - \qty(\dfrac{ku_1}{g+ku_1})\qty(\dfrac{ku_2-g}{k^2}) - \qty(\dfrac{ku_1}{g+ku_1})\qty(\dfrac{g+ku_1}{k^2}) > 0.\]
Hence,
\begin{align*}
	H & > \dfrac{u_1}{g+ku_1}\qty(\dfrac{ku_2-g}{k}) + \dfrac{u_1}{k} \\
	  & = \dfrac{ku_1u_2 +ku_1^2}{k(g+ku_1)}                          \\
	  & = \dfrac{ku_1(u_1+u_2)}{k(g+ku_1)}                            \\
	  & = \dfrac{u_1+u_2}{k}\qty(1 - \dfrac{g}{g+ku_1}).
\end{align*}
However, for $t \neq 0$ and $x_A = 0$, we must have $x_B < 0$.
Denote $t_0 \neq 0$ as the time when $x_A = 0$, so
\begin{align*}
	\dfrac{gt_0}{k} & = \qty(1-e^{-kt_0})\qty(\dfrac{g+ku_1}{k^2}) \\
	t_0             & = \qty(1-e^{-kt_0})\qty(\dfrac{g+ku_1}{gk})  \\
	1 - e^{-kt_0}   & = \dfrac{g+ku_1}{gkt_0}.
\end{align*}
Substituting this into $x_B$, we get:
\begin{align*}
	H & < \qty(1-e^{-kt_0})\qty(\dfrac{ku_2-g}{k^2}) + \qty(1-e^{-kt_0})\qty(\dfrac{g+ku_1}{k^2}) \\
	  & < \qty(\dfrac{ku_2-g}{k^2}) + \qty(\dfrac{g+ku_1}{k^2})                                   \\
	  & = \dfrac{u_1 +u_2}{k}.
\end{align*}
Combining these, we get
\[\boxed{
	  \dfrac{u_1+u_2}{k}\qty(1 - \dfrac{g}{g+ku_1}) < H < \dfrac{u_1+u_2}{k},
}\]
and we are done.
\end{document}
